\chapter{The setting and the main results}
\label{ch:prelims}

We first fix the geometric and probabilistic objects used throughout the blueprint. We then
state the complete classification, its quantitative two-value special case, and the resulting
quasisymmetry statement for visual boundaries. The proofs are assembled in
\cref{ch:gw-simple} and \cref{ch:trichotomy} from the matching theorems of
\cref{ch:iid} and \cref{ch:engine}.

\section{Galton--Watson trees and quasi-isometry}

\begin{definition}[Offspring law and the Galton--Watson tree]
\label{def:gw}
\lean{BranchingProcess.Offspring, BranchingProcess.treeLaw}
\leanok
An offspring law $\theta$ is a probability measure on $\bbN_0$ with finite support. The
Galton--Watson tree $\cT$ with offspring law $\theta$ is the genealogy cut out by a field of
independent offspring counts with law $\theta$, viewed as a subtree of the ambient $N$-ary
tree and carrying the graph metric of its edges.
\end{definition}

\begin{definition}[Quasi-isometry]
\label{def:qi}
\lean{BranchingProcess.IsQIWith, BranchingProcess.QuasiIsometric}
\leanok
A map $\psi\colon(Y_1,d_1)\to(Y_2,d_2)$ is an \emph{$(A,B,C)$-quasi-isometry} if
\[
 A^{-1}d_1(x,y)-B\leq d_2(\psi(x),\psi(y))\leq A\,d_1(x,y)+B
 \qquad(x,y\in Y_1),
\]
and every $y\in Y_2$ lies at distance at most $C$ from $\psi(Y_1)$, where $A\geq1$ and
$B,C\geq0$. Two spaces are quasi-isometric, written $Y_1\simeq Y_2$, when such a map exists.
\end{definition}

\begin{definition}[Quasisymmetry]
\label{def:qs}
Let $\eta\colon[0,\infty)\to[0,\infty)$ be a homeomorphism. A homeomorphism
$\psi\colon(Y_1,d_1)\to(Y_2,d_2)$ is \emph{$\eta$-quasisymmetric} if, for all distinct
$x,y,z\in Y_1$ and all $t>0$,
\[
 d_1(x,z)\leq t\,d_1(y,z)
 \quad\Longrightarrow\quad
 d_2(\psi(x),\psi(z))\leq\eta(t)\,d_2(\psi(y),\psi(z)).
\]
\end{definition}

\begin{definition}[Branching semigroup]
\label{def:branching-semigroup}
\lean{GraphMarkovMatching.shiftedSupportSemigroup}
\leanok
Suppose an offspring law $\theta$ satisfies $\theta_0=0$. Its \emph{branching semigroup} is
the additive submonoid
\[
 \Lambda_\theta
 \coloneqq\bigl\langle k-1 \colon k\geq1\text{ and }\theta_k>0\bigr\rangle
 \subseteq\bbN_0.
\]
\end{definition}

\section{The classification}

\begin{theorem}[Complete quasi-isometry classification]
\label{thm:trichotomy}
\lean{ChainClasses.offspring_classification_ae_iff}
\leanok
\uses{def:gw, def:qi, def:branching-semigroup}
All finite Galton--Watson realisations form the quasi-isometry class $(\mathrm{Fin})$, and no
finite tree is quasi-isometric to an infinite tree.

Let $\theta$ and $\theta'$ be finitely supported offspring laws, each either supercritical or
with $\theta_1=1$, and let $\cT,\cT'$ be independent Galton--Watson trees with these offspring
laws, conditioned on having infinite diameter. Every such law belongs to exactly one of the
following classes:
\begin{itemize}
 \item[(R)] $\theta_1=1$, the \emph{ray} class, in which $\cT$ is the ray.
 \item[(F)] $\theta_0=\theta_1=0$, the \emph{full tree} class, in which $\cT$ is
 quasi-isometric to the binary tree.
 \item[$(\mathrm{C}_\Lambda)$] $\theta_0=0<\theta_1<1$ and $\Lambda_\theta=\Lambda$, the
 \emph{chain} class indexed by the possible branching semigroups $\Lambda$.
 \item[(B)] $\theta_0>0$, the \emph{bushy} class.
\end{itemize}
Two independent realisations in the same class are almost surely quasi-isometric. If they
belong to different classes, then they are almost surely not quasi-isometric.
\end{theorem}

\begin{theorem}[Universality in the two-value family]
\label{thm:twovalue}
\lean{ChainClasses.twovalue_ae_tree_family}
\leanok
\uses{def:gw, def:qi}
Let $\theta$ satisfy $\theta_1+\theta_2=1$, and let $\cT(\omega),\cT(\omega')$ be two
independent Galton--Watson trees with offspring law $\theta$. Then almost surely
$\cT(\omega)\simeq\cT(\omega')$. If moreover $0<\theta_1<1$, there are constants
$D_0=D_0(\theta_1)$ and $C=C(\theta_1)$ such that for every integer $D\geq D_0$,
\[
 \bbP\bigl(\text{there is no }(D^2+3)\text{-quasi-isometry }\cT(\omega)\to\cT(\omega')\bigr)
 \leq C\,\theta_1^{D(D-5/2)}.
\]
\end{theorem}

\begin{corollary}[Quasisymmetry of the boundaries]
\label{thm:boundary}
\uses{thm:trichotomy, def:qs}
Let $\theta$ and $\theta'$ satisfy the hypotheses of \cref{thm:trichotomy}, suppose neither
belongs to the ray class, and suppose they belong to the same class. Let $\cT(\omega)$ and
$\cT'(\omega')$ be independent Galton--Watson trees with these offspring laws, conditioned on
having infinite diameter. Almost surely their visual boundaries are compact, perfect, totally
disconnected metric spaces, and there is a quasisymmetric homeomorphism
\[
 (\partial\cT(\omega),d_{\mathrm{vis}})
 \longrightarrow(\partial\cT'(\omega'),d_{\mathrm{vis}}).
\]
\end{corollary}

\begin{proof}
\uses{thm:trichotomy}
By \cref{thm:trichotomy}, the two tree realisations are almost surely quasi-isometric. Since the
trees are locally finite, they are proper hyperbolic spaces. The theorem of Buyalo and Schroeder
therefore extends the quasi-isometry to a quasisymmetric homeomorphism of their visual
boundaries. This final boundary-extension step is not formalised in Lean.
\end{proof}
